% Partie : logique intuitionniste
% Chapitre : propositions comme types
% Section : types

\documentclass[../../../include/open-logic-section]{subfiles}

\begin{document}

\olfileid[fr]{int}{pty}{typ}

\olsection{Types}

Un terme de preuve tel que $(MN)$ ne détermine pas, à lui seul,
les formules de la preuve qu'il encode. Les variables libres,
c'est-à-dire les termes associés aux axiomes, ne précisent pas
ces formules. Un contexte~$\Gamma$ fournit cette information.
Le type du terme, s'il existe, est alors déterminé de manière
unique. Le terme fixe aussi les règles de la dérivation dans
la présentation à contexte commun donnée ci-dessous.
Dans la présentation \Log{tN2i}, les hypothèses inutilisées
peuvent toutefois être réparties de plusieurs manières entre
les contextes des prémisses.

\begin{defn}
Un \emph{contexte} pour un terme de preuve~$N$ est un ensemble
fini $\Gamma$ de couples $x : !A$, qui attribue au plus un type
à chaque variable et exactement un type à chaque variable
libre de~$N$.
\end{defn}

Le contexte peut aussi attribuer des types à des variables
qui ne sont pas libres dans~$N$.
Par exemple, $\Gamma = \{x:!A, y:!B, z:!C\}$ est un contexte
pour $\pair{x}{y}$.
Nous renommons au besoin les variables liées pour les choisir
distinctes des variables du contexte.

\begin{defn}\ollabel{defn:type}
Soit $\Gamma$ un contexte pour~$N$.
Le \emph{type} de~$N$ relativement à~$\Gamma$, lorsqu'il existe,
est défini inductivement par les clauses suivantes.
Les sous-termes sont typés dans le même contexte,
sauf pour les extensions explicitement indiquées :
\begin{enumerate}
\item La variable $x$ a le type $!A$ si et seulement si
  $x:!A \in \Gamma$.
\item Si $N$ a le type~$!B$ relativement à $x:!A, \Gamma$,
  alors $\lambd[\typeof{x}{!A}][N]$ a le type $!A \lif !B$.
\item Si $N$ a le type $!A \lif !B$ et $M$ le type~$!A$,
  alors $(NM)$ a le type~$!B$.
\item Si $N$ a le type $!A$ et $M$ le type~$!B$, alors
  $\pair{N}{M}$ a le type $!A \land !B$.
\item Si $N$ a le type~$!A_1 \land !A_2$, alors
  $\proj{i}{N}$ a le type~$!A_i$.
\item Si $N$ a le type $!A$, alors $\inj[!B]{i}{N}$
  a le type $!A \lor !B$ pour $i=1$, et le type
  $!B \lor !A$ pour $i=2$.
\item Si $M$ a le type $!A \lor !B$, si $N_1$ a le type
  $!C$ relativement à $x_1:!A, \Gamma$, et si $N_2$
  a le type~$!C$ relativement à $x_2:!B, \Gamma$, alors
  $\dcase{M}{x_1}{N_1}{x_2}{N_2}$ a le type~$!C$.
\item Si $N$ a le type~$\lfalse$, alors $\abort{!A}{N}$
  a le type~$!A$.
\end{enumerate}
\end{defn}

Si les conditions d'une clause ne sont pas satisfaites,
elle n'attribue aucun type au terme.

\begin{prop}
Tout terme de preuve~$N$ a au plus un type relativement
à un contexte~$\Gamma$ pour~$N$.
\end{prop}

\begin{proof}
On raisonne par récurrence sur~$N$.
Le contexte détermine le type d'une variable.
Pour chaque autre constructeur, les types de ses sous-termes,
s'ils existent, sont uniques par l'hypothèse de récurrence.
Avec les annotations du constructeur, ils déterminent
au plus un type pour le terme composé.
Dans une application, les types du domaine et de l'argument
doivent coïncider ; dans une analyse par cas, les deux branches
doivent avoir le même type. Si ces conditions échouent,
le terme n'a pas de type.
\end{proof}

Un terme bien typé possède donc exactement un type dans
le contexte fixé. Nous écrivons $\Gamma \Sequent N : !A$
lorsque $N$ a le type~$!A$ relativement à~$\Gamma$.
Les clauses de la \olref{defn:type} reprennent les règles
de \Log{tN2i} avec un contexte de base~$\Gamma$ commun.
Les règles à lieurs étendent ce contexte dans leurs prémisses.
Elles forment le calcul \Log{tN3i}, présenté dans le
\olref{tab:tN3ip}.

\subfile{rules-tN3}

Les termes de preuve sont les termes d'une variante du
$\lambd$-calcul : le \emph{$\lambd$-calcul typé avec produits,
sommes et type vide}. Le $\lambd$-calcul non typé traditionnel
ne comporte que des variables, des
\emph{$\lambd$-abstractions}~$\lambd[x][N]$ et des applications
$(MN)$. Il ne porte pas l'annotation~$!A$ de l'abstraction
typée $\lambd[\typeof{x}{!A}][N]$.
Le $\lambd$-calcul non typé est un modèle de calcul
Turing-complet : toute fonction partielle calculable
peut y être représentée par un terme.
La variante munie de termes de preuve et de règles de typage
est un modèle plus restreint, dont les principes de calcul
restent les mêmes. Seuls les termes \emph{bien typés},
c'est-à-dire ceux qui possèdent un type relativement
à un contexte~$\Gamma$, y sont admis.
Par exemple, $\lambd[x][(xx)]$ est un terme du lambda-calcul
non typé, mais $\lambd[\typeof{x}{!A}][(xx)]$ n'est bien typé
pour aucune formule $!A$.
Pour typer $(xx)$, la première occurrence de $x$
devrait avoir le type $!A \lif !B$ et la seconde le type~$!A$.
Elles ont pourtant la même variable : ces types devraient
coïncider, ce qui est impossible puisque $!A$ est une
sous-formule propre de~$!A \lif !B$.

Intuitivement, on peut interpréter un type comme la sorte
d'objet désigné par un terme.
Un terme de type $!A \land !B$ désigne un couple dont
les composantes ont respectivement les types~$!A$ et~$!B$,
c'est-à-dire un élément du produit~$!A \times !B$.
Un terme de type $!A \lif !B$ représente une fonction
des objets de type $!A$ vers ceux de type~$!B$.
Un terme de type $!A \lor !B$ indique un objet de type
$!A$ ou de type~$!B$, avec l'information sur le côté choisi.
C'est le principe de l'union disjointe de $!A$ et $!B$,
définie par $\{\tuple{i,x} : i = 1 \text{ ou } 2,
x \in !A \text{ si } i = 1, x \in !B \text{ si } i=2\}$.
Le type $\lfalse$ est interprété par l'ensemble vide :
un terme de ce type ne peut donc désigner aucun objet.
La déduction naturelle intuitionniste, y compris \Log{tN3i},
est cohérente ; elle ne prouve pas~$\lfalse$ sans hypothèse
ouverte. Il n'existe donc pas de terme de preuve clos,
c'est-à-dire sans variable libre, de type~$\lfalse$.

Les opérateurs du lambda-calcul typé produisent des termes
qui désignent les objets attendus, à condition que leurs
arguments aient les types requis.
Ainsi, « si $N_1 : !A$ et $N_2 : !B$, alors
$\pair{N_1}{N_2} : !A \land !B$ » signifie que
si $N_1$ désigne un objet de type~$!A$ et $N_2$ un objet
de type~$!B$, leur couple $\pair{N_1}{N_2}$ désigne
un objet de type $!A \land !B$.

\end{document}
